\documentclass[runningheads]{llncs}
\usepackage[T1]{fontenc}
\usepackage{graphicx}
\usepackage{booktabs}
\usepackage[misc]{ifsym}
\newcommand{\corr}{(\Letter)}
% N.B.: do not change anything above this line. If you require additional packages,please load them directly after this line.
\usepackage{mwe}
% N.B.: you may delete the preceding line. It is used to display an example image in this template.

\begin{document}

\title{Recent Advances in Underwater Basket Weaving Under the Extreme Pressure of the Mariana Trench}

\titlerunning{Underwater Basket Weaving Under Extreme Pressure}
% If the full title of your paper is short enough to also fit in the running head,you can omit the abbreviated paper title here. You can check as follows: if you comment out the \titlerunning line,something will appear in the header of all odd-numbered pages of your PDF from page 3 onward. This something is either the full title (in which case all is well),or the error message "Title Suppressed Due to Excessive Length". If this error message appears,you're going to want to provide an abbreviated title within the \titlerunning command,because if you won't do it,Springer will do it for you.

%N.B.: Author information (both in the \author{} and \authorrunning{} command) should only be present in the Camera-Ready Version of your paper. The version that you initially submit for review,ought to be double-blind. So,when initially submitting your paper,use:
%\author{Author information scrubbed for double-blind reviewing}
\author{Andr\'e Lauren Benjamin\inst{1} \and
Calvin Cordozar Broadus Jr.\inst{2,3} \corr \and
Antwan Andr\'e Patton\inst{1}\orcidID{0000-1111-2222-3333}}
% You may leave out the orcidID information,if you want to.
% Use \corr to indicate the corresponding author. Note the spacing around the \corr command. Only one author can be the corresponding author.

%N.B.: comment out the \authorrunning{} command for the double-blind version of your paper submitted for review. Later,if your paper is accepted,use the command for the Camera-Ready Version.
\authorrunning{A.L. Benjamin et al.}
% First names are abbreviated in the running head.
% If there is one author,write 'A.L. Benjamin'.
% If there are two authors,write 'A.L. Benjamin and C.C. Broadus Jr.'
% If there are more than two authors,'[...] et al.' is used.

\institute{Fictional Southern University,Savannah GA 31404,USA \email{\{a.l.benjamin,a.a.patton\}@fsu.fake}
\and
Fictional West Coast University,Long Beach CA 90840,USA \email{ccb@fwcu.fake}
\and
Secondary European Affiliation,Tiergartenstr. 17,69121 Heidelberg,Germany
\email{lncs@springer.com}}

\maketitle              % typeset the header of the contribution

\begin{abstract}
The abstract should briefly summarize the contents of the paper in
150--250 words.

\keywords{First keyword  \and Second keyword \and Another keyword.}
\end{abstract}

\section{Introduction}
\subsection{Main Contributions}
Please note that the first paragraph of a section or subsection is
not indented. The first paragraph that follows a table,figure,
equation etc. does not need an indent,either.

Subsequent paragraphs,however,are indented.

\section{Related Work}
\subsection{Basket Weaving}
\subsubsection{Underwater Basket Weaving} Only two levels of
headings should be numbered. Lower level headings remain unnumbered;
they are formatted as run-in headings.

\paragraph{Underwater Basket Weaving Under Difficult Circumstances}
The contribution should contain no more than four levels of
headings. Table~\ref{tab1} gives a summary of all heading levels.

\begin{table}[t]
\caption{Table captions should be placed above the
tables.}\label{tab1}
\begin{tabular}{lll}
\toprule
Heading level &  Example & Font size and style\\
\midrule
Title (centered) &  {\Large\bfseries Lecture Notes} & 14 point,bold\\
1st-level heading &  {\large\bfseries 1 Introduction} & 12 point,bold\\
2nd-level heading & {\bfseries 2.1 Printing Area} & 10 point,bold\\
3rd-level heading & {\bfseries Run-in Heading in Bold.} Text follows & 10 point,bold\\
4th-level heading & {\itshape Lowest Level Heading.} Text follows & 10 point,italic\\
\bottomrule
\end{tabular}
\end{table}

\section{Recent Advances from the Mariana Trench}

Displayed equations are centered and set on a separate
line.
\begin{equation}
x + y = z
\end{equation}
Please try to avoid rasterized images for line-art diagrams and
schemas. Whenever possible,use vector graphics instead (see
Fig.\@ \ref{fig1}).

\begin{figure}[t]
\includegraphics[width=\textwidth]{example-image-duck}
\caption{A figure caption is always placed below the illustration.
Please note that short captions are centered,while long ones are
justified by the macro package automatically.} \label{fig1}
\end{figure}
%Note that the \includegraphics command is commented out in this example figure,so that we can compile this document without sending sample figures around.

\begin{theorem}
This is a sample theorem. The run-in heading is set in bold,while
the following text appears in italics. Definitions,lemmas,
propositions,and corollaries are styled the same way.
\end{theorem}
%
% the environments 'definition','lemma','proposition','corollary',
% 'remark',and 'example' are defined in the LLNCS documentclass as well.
%
\begin{proof}
Proofs,examples,and remarks have the initial word in italics,
while the following text appears in normal font.
\end{proof}
For citations of references,we prefer the use of square brackets
and consecutive numbers. Citations using labels or the author/year
convention are also acceptable. The following bibliography provides
a sample reference list with entries for journal
articles~\cite{ref_article1},an LNCS chapter~\cite{ref_lncs1},a
book~\cite{ref_book1},proceedings without editors~\cite{ref_proc1},
and a homepage~\cite{ref_url1}. Multiple citations are grouped
\cite{ref_article1,ref_lncs1,ref_book1},
\cite{ref_article1,ref_book1,ref_proc1,ref_url1}.

\section{Experiments}
\subsection{Experimental Setup}
\subsection{Experimental Results}

\section{Discussion}

\section{Conclusion}

Of course,authors have complete freedom on how they choose to structure their paper. Section headers from Introduction up to and including Conclusions are completely up to the discretion of the authors; use whichever structure you see fit. Title,Abstract,the credits environment,and References,however,are mandatory.

\begin{credits}
\subsubsection{\ackname} A bold run-in heading in small font size at the end of the paper is
used for general acknowledgments,for example: This study was funded
by X (grant number Y).

\subsubsection{\discintname}
It is now necessary to declare any competing interests or to specifically
state that the authors have no competing interests. Please place the
statement with a bold run-in heading in small font size beneath the
(optional) acknowledgments,
for example: The authors have no competing interests to declare that are
relevant to the content of this article. Or: Author A has received research
grants from Company W. Author B has received a speaker honorarium from
Company X and owns stock in Company Y. Author C is a member of committee Z.
\end{credits}
%
% ---- Bibliography ----
%
% BibTeX users should specify bibliography style 'splncs04'.
% References will then be sorted and formatted in the correct style.
%
% \bibliographystyle{splncs04}
% \bibliography{mybibliography}
%% Note that this preceding line implies that you store your BibTeX references in a file called 'mybibliography.bib'. If you instead store your references in a file with a different name,for instance 'references.bib',the preceding line should read '\bibliography{references}'. Whatever you do,DO NOT put the file name extension .bib inside the \bibliography command; this will trip up LaTeX compilers. 
%
% If you do not want to use BibTeX,you can also type up the bibliography exactly as you see fit,using the following structure:
\begin{thebibliography}{8}
% Note that this number 8 reserves an amount of space (equal to the natural width of the given number) for the label of your references; if you have more than 9 references,you will want to change this number to 18. If you have more than 19 references,this number is best changed to 88. If you have more than 99 references,I salute you.
\bibitem{ref_article1}
Author,F.: Article title. Journal \textbf{2}(5),99--110 (2016)

\bibitem{ref_lncs1}
Author,F.,Author,S.: Title of a proceedings paper. In: Editor,
F.,Editor,S. (eds.) CONFERENCE 2016,LNCS,vol. 9999,pp. 1--13.
Springer,Heidelberg (2016). \doi{10.10007/1234567890}

\bibitem{ref_book1}
Author,F.,Author,S.,Author,T.: Book title. 2nd edn. Publisher,
Location (1999)

\bibitem{ref_proc1}
Author,A.-B.: Contribution title. In: 9th International Proceedings
on Proceedings,pp. 1--2. Publisher,Location (2010)

\bibitem{ref_url1}
LNCS Homepage,\url{http://www.springer.com/lncs},last accessed 2023/10/25
\end{thebibliography}
\end{document}
